\section{The key formula}
\label{section-key}

\noindent
Let $(S, \delta)$ be as in Situation \ref{situation-setup}.
Let $X$ be locally of finite type over $S$. Assume
$X$ is integral and $\dim_\delta(X) = n$.
Let $\mathcal{L}$ and $\mathcal{N}$ be invertible sheaves on $X$.
Let $s$ be a nonzero meromorphic section of $\mathcal{L}$ and
let $t$ be a nonzero meromorphic section of $\mathcal{N}$.
Let $Z_i \subset X$, $i \in I$ be a locally finite set of irreducible
closed subsets of codimension $1$ with the following property:
If $Z \not \in \{Z_i\}$ with generic point $\xi$, then $s$ is a generator
for $\mathcal{L}_\xi$ and $t$ is a generator for $\mathcal{N}_\xi$.
Such a set exists by
Divisors, Lemma \ref{divisors-lemma-divisor-meromorphic-locally-finite}.
Then
$$
\text{div}_\mathcal{L}(s) = \sum \text{ord}_{Z_i, \mathcal{L}}(s) [Z_i]
$$
and similarly
$$
\text{div}_\mathcal{N}(t) = \sum \text{ord}_{Z_i, \mathcal{N}}(t) [Z_i]
$$
Unwinding the definitions more, we pick for each $i$ generators
$s_i \in \mathcal{L}_{\xi_i}$ and $t_i \in \mathcal{N}_{\xi_i}$
where $\xi_i$ is the generic point of $Z_i$. Then we can write
$$
s = f_i s_i
\quad\text{and}\quad
t = g_i t_i
$$
Set $B_i = \mathcal{O}_{X, \xi_i}$. Then by definition
$$
\text{ord}_{Z_i, \mathcal{L}}(s) = \text{ord}_{B_i}(f_i)
\quad\text{and}\quad
\text{ord}_{Z_i, \mathcal{N}}(t) = \text{ord}_{B_i}(g_i)
$$
Since $t_i$ is a generator of $\mathcal{N}_{\xi_i}$ we see that
its image in the fibre $\mathcal{N}_{\xi_i} \otimes \kappa(\xi_i)$
is a nonzero meromorphic section of $\mathcal{N}|_{Z_i}$. We will denote
this image $t_i|_{Z_i}$. From our definitions it follows that
$$
c_1(\mathcal{N}) \cap \text{div}_\mathcal{L}(s) =
\sum \text{ord}_{B_i}(f_i)
(Z_i \to X)_*\text{div}_{\mathcal{N}|_{Z_i}}(t_i|_{Z_i})
$$
and similarly
$$
c_1(\mathcal{L}) \cap \text{div}_\mathcal{N}(t) =
\sum \text{ord}_{B_i}(g_i)
(Z_i \to X)_*\text{div}_{\mathcal{L}|_{Z_i}}(s_i|_{Z_i})
$$
in $\CH_{n - 2}(X)$. We are going to find a rational equivalence between
these two cycles. To do this we consider the tame symbol
$$
\partial_{B_i}(f_i, g_i) \in \kappa(\xi_i)^*
$$
see Section \ref{section-tame-symbol}.

\begin{lemma}[Key formula]
\label{lemma-key-formula}
In the situation above the cycle
$$
\sum
(Z_i \to X)_*\left(
\text{ord}_{B_i}(f_i) \text{div}_{\mathcal{N}|_{Z_i}}(t_i|_{Z_i}) -
\text{ord}_{B_i}(g_i) \text{div}_{\mathcal{L}|_{Z_i}}(s_i|_{Z_i}) \right)
$$
is equal to the cycle
$$
\sum (Z_i \to X)_*\text{div}(\partial_{B_i}(f_i, g_i))
$$
\end{lemma}

\begin{proof}
First, let us examine what happens if we replace $s_i$ by $us_i$
for some unit $u$ in $B_i$. Then $f_i$ gets replaced by $u^{-1} f_i$.
Thus the first part of the first expression of the lemma is unchanged
and in the second part we add
$$
-\text{ord}_{B_i}(g_i)\text{div}(u|_{Z_i})
$$
(where $u|_{Z_i}$ is the image of $u$ in the residue field) by
Divisors, Lemma \ref{divisors-lemma-divisor-meromorphic-well-defined}
and in the second expression we add
$$
\text{div}(\partial_{B_i}(u^{-1}, g_i))
$$
by bi-linearity of the tame symbol. These terms agree by property
(\ref{item-normalization}) of the tame symbol.

\medskip\noindent
Let $Z \subset X$ be an irreducible closed with $\dim_\delta(Z) = n - 2$.
To show that the coefficients of $Z$ of the two cycles of the lemma
is the same, we may do a replacement $s_i \mapsto us_i$ as in the previous
paragraph. In exactly the same way one shows that we may do a replacement
$t_i \mapsto vt_i$ for some unit $v$ of $B_i$.

\medskip\noindent
Since we are proving the equality of cycles we may argue one coefficient
at a time. Thus we choose an irreducible closed $Z \subset X$
with $\dim_\delta(Z) = n - 2$ and compare coefficients. Let $\xi \in Z$
be the generic point and set $A = \mathcal{O}_{X, \xi}$. This is a Noetherian
local domain of dimension $2$. Choose generators $\sigma$ and $\tau$
for $\mathcal{L}_\xi$ and $\mathcal{N}_\xi$. After shrinking $X$, we may
and do assume $\sigma$ and $\tau$ define trivializations
of the invertible sheaves $\mathcal{L}$ and $\mathcal{N}$ over all of $X$.
Because $Z_i$ is locally
finite after shrinking $X$ we may assume $Z \subset Z_i$ for all $i \in I$
and that $I$ is finite. Then $\xi_i$ corresponds to a prime
$\mathfrak q_i \subset A$ of height $1$.
We may write $s_i = a_i \sigma$ and $t_i = b_i \tau$
for some $a_i$ and $b_i$ units in $A_{\mathfrak q_i}$.
By the remarks above, it suffices to prove the lemma when
$a_i = b_i = 1$ for all $i$.

\medskip\noindent
Assume $a_i = b_i = 1$ for all $i$. Then the first expression of the
lemma is zero, because we choose $\sigma$ and $\tau$ to be trivializing
sections. Write $s = f\sigma$ and $t = g \tau$ with $f$ and $g$ in the
fraction field of $A$. By the previous paragraph we have reduced to the case
$f_i = f$ and $g_i = g$ for all $i$. Moreover, for a height $1$ prime
$\mathfrak q$ of $A$ which is not in $\{\mathfrak q_i\}$ we have
that both $f$ and $g$ are units in $A_\mathfrak q$ (by our choice of
the family $\{Z_i\}$ in the discussion preceding the lemma). Thus
the coefficient of $Z$ in the second expression of the lemma is
$$
\sum\nolimits_i \text{ord}_{A/\mathfrak q_i}(\partial_{B_i}(f, g))
$$
which is zero by the key Lemma \ref{lemma-milnor-gersten-low-degree}.
\end{proof}

\begin{remark}
\label{remark-higher-chow-pointwise}
Let $(S, \delta)$ be as in Situation \ref{situation-setup}.
Let $X$ be locally of finite type over $S$. Let $k \in \mathbf{Z}$.
We claim that there is a complex
$$
\bigoplus\nolimits_{\delta(x) = k + 2}' K_2^M(\kappa(x))
\xrightarrow{\partial}
\bigoplus\nolimits_{\delta(x) = k + 1}' K_1^M(\kappa(x))
\xrightarrow{\partial}
\bigoplus\nolimits_{\delta(x) = k}' K_0^M(\kappa(x))
$$
Here we use notation and conventions introduced in
Remark \ref{remark-chow-group-pointwise} and in addition
\begin{enumerate}
\item $K_2^M(\kappa(x))$ is the degree $2$ part of
the Milnor K-theory of the residue field $\kappa(x)$ of the point
$x \in X$ (see Remark \ref{remark-gersten-complex-milnor}) which
is the quotient of $\kappa(x)^* \otimes_\mathbf{Z} \kappa(x)^*$
by the subgroup generated by elements of the form
$\lambda \otimes (1 - \lambda)$ for
$\lambda \in \kappa(x) \setminus \{0, 1\}$, and
\item the first differential $\partial$ is defined as follows:
given an element $\xi = \sum_x \alpha_x$ in the first term
we set
$$
\partial(\xi) = \sum\nolimits_{x \leadsto x',\ \delta(x') = k + 1}
\partial_{\mathcal{O}_{W_x, x'}}(\alpha_x)
$$
where
$\partial_{\mathcal{O}_{W_x, x'}} : K_2^M(\kappa(x)) \to K_1^M(\kappa(x))$
is the tame symbol constructed in Section \ref{section-tame-symbol}.
\end{enumerate}
We claim that we get a complex, i.e., that $\partial \circ \partial = 0$.
To see this it suffices to take an element $\xi$ as above and a point
$x'' \in X$ with $\delta(x'') = k$ and check that the coefficient of
$x''$ in the element $\partial(\partial(\xi))$ is zero.
Because $\xi = \sum \alpha_x$ is a locally finite sum, we
may in fact assume by additivity that $\xi = \alpha_x$ for
some $x \in X$ with $\delta(x) = k + 2$ and $\alpha_x \in K_2^M(\kappa(x))$.
By linearity again we may assume that $\alpha_x = f \otimes g$ for
some $f, g \in \kappa(x)^*$. Denote $W \subset X$ the integral closed
subscheme with generic point $x$. If $x'' \not \in W$, then it is
immediately clear that the coefficient of $x$ in $\partial(\partial(\xi))$
is zero. If $x'' \in W$, then we see that the coefficient of $x''$
in $\partial(\partial(x))$ is equal to
$$
\sum\nolimits_{x \leadsto x' \leadsto x'',\ \delta(x') = k + 1}
\text{ord}_{\mathcal{O}_{\overline{\{x'\}}, x''}}(
\partial_{\mathcal{O}_{W, x'}}(f, g))
$$
The key algebraic Lemma \ref{lemma-milnor-gersten-low-degree}
says exactly that this is zero.
\end{remark}

\begin{remark}
\label{remark-higher-chow}
Let $(S, \delta)$ be as in Situation \ref{situation-setup}.
Let $X$ be locally of finite type over $S$. Let $k \in \mathbf{Z}$.
The complex in Remark \ref{remark-higher-chow-pointwise} and the
presentation of $\CH_k(X)$ in Remark \ref{remark-chow-group-pointwise}
suggests that we can define a first higher Chow group
$$
\CH^M_k(X, 1) =
H_1(\text{the complex of Remark \ref{remark-higher-chow-pointwise}})
$$
We use the supscript ${}^M$ to distinguish our notation from the
higher chow groups defined in the literature, e.g., in the papers
by Spencer Bloch (\cite{Bloch} and \cite{Bloch-moving}).
Let $U \subset X$ be open with complement $Y \subset X$ (viewed as reduced
closed subscheme). Then we find a split short exact sequence
$$
0 \to
\bigoplus\nolimits_{y \in Y, \delta(y) = k + i}' K_i^M(\kappa(y)) \to
\bigoplus\nolimits_{x \in X, \delta(x) = k + i}' K_i^M(\kappa(x)) \to
\bigoplus\nolimits_{u \in U, \delta(u) = k + i}' K_i^M(\kappa(u)) \to 0
$$
for $i = 2, 1, 0$ compatible with the boundary maps in the complexes
of Remark \ref{remark-higher-chow-pointwise}. Applying the snake lemma
(see Homology, Lemma \ref{homology-lemma-long-exact-sequence-chain})
we obtain a six term exact sequence
$$
\CH^M_k(Y, 1) \to \CH^M_k(X, 1) \to \CH^M_k(U, 1) \to
\CH_k(Y) \to \CH_k(X) \to \CH_k(U) \to 0
$$
extending the canonical exact sequence of Lemma \ref{lemma-restrict-to-open}.
With some work, one may also define flat pullback and proper pushforward
for the first higher chow group $\CH^M_k(X, 1)$. We will return to this
later (insert future reference here).
\end{remark}






